\paragraph{可复核接口：\texorpdfstring{$\mathrm{Aut}(Q_6)$}{Aut(Q6)} 几何稳定子（唯一 \texorpdfstring{$\ZZ_2$}{Z2}）与 \texorpdfstring{$\delta=34$}{δ=34} 掩码指纹}
除“纤维内任意置换”的 fold-gauge 群外，正文在小节 \ref{subsec:bdry-tower-zeck-gut} 进一步引入了一个几何层的稳定子：
将 \(F_6(\omega)=\mathrm{Fold}^{\mathrm{bin}}_6(\mathrm{int}_6(\omega))\) 视为超立方体图 \(Q_6\) 顶点集上的标签函数，并在 \(\mathrm{Aut}(Q_6)\) 内筛选与之交换者，得到唯一残存的非平凡 \(\ZZ_2\) 对称（定理 \ref{thm:foldbin6-geo-stabilizer-z2}），其在二进制编号上携带掩码整数 \(34=100010_2\) 并在半立方体上诱导 \(\pm 34\) 偏移（推论 \ref{cor:foldbin6-geo-mask-34}）。
\begin{itemize}
  \item \textbf{生成脚本：}\texttt{\detokenize{scripts/exp_foldbin6_geo_stabilizer.py}}。
  \item \textbf{LaTeX 生成物：}\texttt{\detokenize{sections/generated/eq_foldbin6_geo_stabilizer.tex}}。
  \item \textbf{审计 JSON：}\texttt{\detokenize{artifacts/export/foldbin6_geo_stabilizer.json}}（含 \(|\mathrm{Stab}_{\mathrm{geo}}|=2\) 的穷举证书、生成元的 \((\mathrm{mask},\mathrm{perm})\) 参数化、固定点/支撑计数、以及二进制编号差值直方图）。
  \item \textbf{几何审计接口（\(\sigma_{\mathrm{geo}}\) 在纤维内的轨道电荷分裂）。}
  \texttt{\detokenize{scripts/exp_foldbin6_geo_orbit_charge.py}} 生成表
  \texttt{\detokenize{sections/generated/tab_foldbin6_geo_orbit_charge.tex}}
  与证书 \texttt{\detokenize{artifacts/export/foldbin6_geo_orbit_charge.json}}，
  给出对每个 \(x\in X_6\) 的 \((f_x,p_x)\)（不动点数/二循环数）与七类直方图，可直接核验推论 \ref{cor:terminal-foldbin6-geo-orbit-charge}。
  \item \textbf{动力学审计接口（强可并性失败的显式反例）。}
  \path{scripts/exp_foldbin6_strong_lumpability_counterexample.py} 生成片段
  \texttt{\detokenize{sections/generated/eq_foldbin6_strong_lumpability_counterexample.tex}}
  与证书 \texttt{\detokenize{artifacts/export/foldbin6_strong_lumpability_counterexample.json}}，
  给出同纤维两点的一步邻域稳定类型计数分布差异，可直接核验“Foldbin6 纤维分割不具强可并（lumpability）”的断言。
  \item \textbf{几何审计接口（纤维内部最小汉明距离三值律）。}
  \path{scripts/exp_foldbin6_fiber_hamming_min_distance.py} 生成表
  \path{sections/generated/tab_foldbin6_fiber_hamming_min_distance.tex}
  与计数直方图片段 \path{sections/generated/eq_foldbin6_fiber_hamming_min_distance_value_counts.tex}
  与证书 \path{artifacts/export/foldbin6_fiber_hamming_min_distance.json}，
  给出逐 \(w\in X_6\) 的 \(\delta_{\min}(w)\) 与预像清单，可直接核验定理 \ref{thm:terminal-foldbin6-fiber-hamming-three-valued} 的 \(\{2,3,5\}\) 三值律。
  \item \textbf{几何审计接口（类型邻接图 \texorpdfstring{$\Gamma_6$}{Gamma6} 的刚性）。}
  \texttt{\detokenize{scripts/exp_window6_type_adjacency_graph_rigidity.py}} 生成片段
  \texttt{\detokenize{sections/generated/eq_window6_type_adjacency_graph_rigidity.tex}}
  与证书 \texttt{\detokenize{artifacts/export/window6_type_adjacency_graph_rigidity.json}}，
  给出 \(|V|,|E|\)、度直方图、连通性、直径、单谱性与 \(\mathrm{Aut}(\Gamma_6)=\{1\}\) 的可复核结论（命题 \ref{prop:terminal-gamma6-rigidity} 与推论 \ref{cor:terminal-gamma6-diameter-simple-spectrum}）。
\end{itemize}

\begin{corollary}[6D 可观测判定的有限完备性：由 \texorpdfstring{$(\pi_6,P_6,\sigma_{\mathrm{geo}})$}{(pi6,P6,sigma_geo)} 给出的可判定同构模板]\label{cor:terminal-window6-finite-completeness-template}
在 window-\(6\) 的审计口径下，令
\[
\mathfrak{A}_6:=(X_6,\pi_6,P_6,\langle\sigma_{\mathrm{geo}}\rangle),
\]
其中 \(\pi_6\) 与 \(P_6\) 由定理 \ref{thm:terminal-foldbin6-pushforward-markov} 给出，\(\langle\sigma_{\mathrm{geo}}\rangle\) 由定理 \ref{thm:foldbin6-geo-stabilizer-z2} 给出。
则对任何另一套在 \(m=6\) 层级输出同构数据的实现（例如来自任意满足命题 \ref{prop:terminal-cut-project-to-fold-hist} 的 6D cut-and-project 实例），
若其在某个标号下给出同一三元组 \((\pi_6,P_6,\langle\sigma_{\mathrm{geo}}\rangle)\)，则该标号由加权邻接刚性推论 \ref{cor:terminal-gamma6-multiplicity-rigidity} 唯一确定；
从而“是否同构”为有限可判定问题，并且任一偏离都可由有限证书直接证伪。
\end{corollary}
\leanverified{paper\_terminal\_window6\_finite\_completeness\_template}
\begin{proof}
给定 \(\pi_6\) 与 \(P_6\) 可恢复推送邻接计数 \(A_6\) 以及相应的加权邻接多重度矩阵 \(A=(A_{uv})\)（见定理 \ref{thm:terminal-foldbin6-pushforward-markov} 的定义式）。
推论 \ref{cor:terminal-gamma6-multiplicity-rigidity} 排除任何保持 \(A\) 的非平凡重标号，从而标号唯一。
几何 \(\ZZ_2\) 稳定子由定理 \ref{thm:foldbin6-geo-stabilizer-z2} 唯一给出。
因此比较两套实现时，只需对有限对象作等式核验即可，结论成立。证毕。
\end{proof}

\begin{remark}[结构定义的可审计唯一性（window-\(6\)）]\label{rem:terminal-window6-structure-defined-uniqueness}
由命题 \ref{prop:terminal-gamma6-rigidity} 与推论 \ref{cor:terminal-gamma6-multiplicity-rigidity}，任何仅依赖于 \(Q_6\) 的边推送数据（等价地：仅依赖于加权邻接多重度矩阵 \(A\)，或由其构造出的 \(P_6\)、edge--flux skeleton 等）的结构，在同构意义下均无非平凡重标号自由度。
因此，凡以“由 \(Q_6\) 局部扰动推送到 \(X_6\)”为函子性约束所定义的附加对象（分块、粗粒化、谱指纹、算术不变量等），都可被视为标签规范且可审计的有限对象；任一声称的替代实现若在这些有限输出上不一致，即可被有限证书直接证伪。
\end{remark}

\begin{remark}[几何刚性消除重标号自由：三因子的规范定位（window-\(6\)）]\label{rem:window6-no-relabeling-factor-localization}
命题 \ref{prop:terminal-gamma6-rigidity} 给出 \(\mathrm{Aut}(\Gamma_6)=\{1\}\)，从而任何保持局部邻接结构的类型重标号均为恒等；
因此由 \(Q_6\) 边推送所定义的 boundary/cyclic 划分、根字典与 \(P_6\) 等输出在标号意义下皆为规范（亦见备注 \ref{rem:terminal-window6-structure-defined-uniqueness}）。
与定理 \ref{thm:foldbin6-geo-stabilizer-z2} 的唯一几何稳定子 \(\langle\sigma_{\mathrm{geo}}\rangle\cong\ZZ_2\) 并置可知：
\(\sigma_{\mathrm{geo}}\) 不能被任何重标号吸收。
故边界词塔给出的 \((1,3,8)\) 维数接口与推论 \ref{cor:bdry-tower-minimal-resonance-triple-sm} 的
\(\mathfrak u(1)\oplus\mathfrak{su}(2)\oplus\mathfrak{su}(3)\) 分解在本审计口径下具有规范定位：
任一声称的“因子互换/层级重新分配”若保持局部邻接与推送核数据，则必与 \(\mathrm{Aut}(\Gamma_6)=\{1\}\) 冲突。
\par
进一步地，由于几何层的稳定子仅为 \(\langle\sigma_{\mathrm{geo}}\rangle\cong\ZZ_2\)（定理 \ref{thm:foldbin6-geo-stabilizer-z2}），
几何侧无法在类型层诱导任何非平凡重标号自由度；
因此凡在 \(m=6\) 的读出层面出现的非平凡对称（例如 boundary sheet-flip 与其中心嵌入，推论 \ref{cor:terminal-delta-parity-rule} 与推论 \ref{cor:fold-bin-boundary-center-z2}），
只能来自 fold-gauge 群 \(\mathcal{G}^{\mathrm{bin}}_6\) 的纤维内置换，而非来自 \(\mathrm{Aut}(Q_6)\) 的几何自同构。
\end{remark}

\paragraph{可复核接口：独立位翻转噪声下的 \texorpdfstring{$\mathrm{Fold}^{\mathrm{bin}}_6$}{Foldbin6} 精确稳定多项式}
终端命题定理 \ref{thm:terminal-foldbin6-bitflip-stability-polynomial} 将 window-\(6\) 的稳定类型读出在 i.i.d.\ 位翻转信道下的稳定概率写成关于翻转率 \(p\) 的闭式多项式；其系数由 \(64\) 微态的有限穷举严格确定：
\begin{itemize}
  \item \textbf{生成脚本：}\texttt{\detokenize{scripts/exp_foldbin6_bitflip_stability_polynomial.py}}。
  \item \textbf{LaTeX 生成物：}\texttt{\detokenize{sections/generated/eq_foldbin6_bitflip_stability_polynomial.tex}}。
  \item \textbf{审计 JSON：}\texttt{\detokenize{artifacts/export/foldbin6_bitflip_stability_polynomial.json}}（含按噪声重量分层的稳定对数与多项式系数）。
\end{itemize}

\paragraph{可复核接口：纤维加权耦合矩阵 \(M\) 的可逆性与 \texorpdfstring{$\det(M)$}{det(M)} 素因子分解}
终端命题定理 \ref{thm:terminal-window6-fiber-edge-coupling-det} 在 \(X_6\) 上以 \(Q_6\) 的边诱导加权耦合矩阵 \(M\in\mathbb{Z}^{21\times 21}\)（纤维间边数），并给出行和可观测性恒等式与 \(\det(M)\neq 0\) 的整因子分解证书：
\begin{itemize}
  \item \textbf{生成脚本：}\texttt{\detokenize{scripts/exp_window6_fiber_edge_coupling_matrix.py}}。
  \item \textbf{LaTeX 生成物：}\texttt{\detokenize{sections/generated/eq_window6_fiber_edge_coupling_det.tex}}。
  \item \textbf{审计 JSON：}\texttt{\detokenize{artifacts/export/window6_fiber_edge_coupling_matrix.json}}（含 \(M\) 的全矩阵、纤维大小与行和核验）。
\end{itemize}

\paragraph{可复核接口：\texorpdfstring{$X_6$}{X6} 的 \(1+8+12\) 字典与 \(12=8+4\) 两档重谱}
“\(1+8+12\)” 与 “\(12=8+4\)” 的两条离散分块结论无需另写脚本：
它们可直接从表 \ref{tab:fold6_bin_uplift_choice_collapse} 的三行词表读出。
第一行给出 \(S_{\mathrm{low}}\) 的 \(9\) 词（其中 \(\texttt{000000}\) 为唯一零词、其余 \(8\) 词构成 \(1+8\) 的轻块），
第二行与第三行合并为 \(12\) 的重块，并按 \(w_6=1\)（\(\Xi=2\)，\(8\) 词）与 \(w_6=0,\ V_6(w)>8\)（\(\Xi=3\)，\(4\) 词）自动分成两档。
由此可直接计算熵权重 \(h(w)=\log\Xi(w)\) 并作为谱隙对齐的分块权重（\S\ref{subsec:group-unification-spectral-alignment}）。
对应的可复核输出与脚本沿用上文关于 \(C_6\) 轨道与 dyadic uplift 证书的生成链（表 \ref{tab:fold6_bin_uplift_choice_collapse} 与
\texttt{\detokenize{artifacts/export/window6_c6_orbit_patisalam_seed.json}}）。

\paragraph{可复核接口：bin-fold 退化谱的双尺度选择律与 \texorpdfstring{$\dim=55$}{dim=55} 统一锚点}
bin-fold 退化谱的双尺度选择律将二进制区间折叠 \(\mathrm{Fold}^{\mathrm{bin}}_m\) 的退化度谱
\(
d^{\mathrm{bin}}_m(w)=|\mathrm{Fold}^{\mathrm{bin}}_m{}^{-1}(w)|
\)
作为独立于 Zeckendorf--GUT 签名与 NAP 选择器的第三条审计通道：
在 \(m=6\) 上由低退化扇区读出 \(12\) 的计数接口，并以 ghost 容量 \(H^{\mathrm{bin}}_6=2^6-|X_6|=43\) 补全得到 \(55\) 的统一维数锚点；
在更高分辨率 \(m=10\) 上，最小退化扇区计数再次精确等于 \(55\)，从而形成两路锁定。
其生成链已脚本化并纳入一键流水线：
\begin{itemize}
  \item \textbf{生成脚本：}\texttt{\detokenize{scripts/exp_fold_bin_degeneracy_spectrum_m6_m12.py}}。
  \item \textbf{LaTeX 生成物：}\texttt{\detokenize{sections/generated/tab_fold_bin_degeneracy_spectrum_m6_m12.tex}}（表 \ref{tab:fold_bin_degeneracy_spectrum_m6_m12}）。
  \item \textbf{审计 JSON：}\texttt{\detokenize{artifacts/export/fold_bin_degeneracy_spectrum_m6_m12.json}}（含各 \(m\) 的直方图、\(d_{\min}\)、最小扇区大小与墙钟时间）。
\end{itemize}

\paragraph{交叉审计接口：\texorpdfstring{$\mathrm{Fold}^{\mathrm{bin}}$}{Foldbin}–\texorpdfstring{$\mathrm{Fold}$}{Fold} 的跨尺度补余恒等式}
接口化结论 \ref{par:gut-foldbin-fold-complement-identity} 给出一条把“低分辨率最刚性扇区补余容量”与“高分辨率最大纤维极值模尺度”直接锁定的交叉审计恒等式；
其核对式为独立的派生输出（不重新枚举谱数据），并已纳入一键流水线：
\begin{itemize}
  \item \textbf{生成脚本：}\texttt{\detokenize{scripts/exp_foldbin_fold_complement_identity_m6_m12.py}}。
  \item \textbf{LaTeX 生成物：}\texttt{\detokenize{sections/generated/eq_foldbin_fold_complement_identity_m6_m12.tex}}。
  \item \textbf{审计 JSON：}\texttt{\detokenize{artifacts/export/foldbin_fold_complement_identity_m6_m12.json}}（含逐 \(m\) 核对项与通过/失败标记）。
\end{itemize}

\paragraph{可复核接口：四步平移 uplift 同构与 \texorpdfstring{$m=10$}{m=10} 的 \(6+6+6+3\) 根型分块}
四步平移 uplift 同构把 \(X_{10}^{\mathrm{bdry}}\) 与 \(X_6\) 规范同构，并以 \(B_3\) 根字典细化为
\(6\) 个短根、\(6\) 个 \(A_2(+)\)、\(6\) 个 \(A_2(-)\) 与 \(3\) 个 Cartan 方向。
对应的 \(21\) 词 uplift 字典与分块清单由脚本自动生成：
\begin{itemize}
  \item \textbf{生成脚本：}\texttt{\detokenize{scripts/exp_boundary_uplift_m10_b3c3_dictionary.py}}。
  \item \textbf{LaTeX 生成物：}\texttt{\detokenize{sections/generated/eq_boundary_uplift_m10_isomorphism.tex}}，
  \texttt{\detokenize{sections/generated/tab_boundary_uplift_m10_dictionary.tex}}。
  \item \textbf{审计 JSON：}\texttt{\detokenize{artifacts/export/boundary_uplift_m10_b3c3_dictionary.json}}（含 \(\iota_n\) 的有限窗双射核验、\(m=10\) 的 21 词字典与 \(A_2(\pm)\) 分块清单）。
\end{itemize}

\begin{remark}[长根的双 \texorpdfstring{$A_2$}{A2} 子系分裂与符号刚性]\label{rem:boundary-m10-longroot-double-a2-sign-rigidity}
在定理 \ref{thm:boundary-m10-b3c3-rigid-four-block-split} 的口径下，
\(m=10\) 的 \(18\) 条根方向严格分解为
\(
6(\mathrm{short})\ \sqcup\ 6(A_2(+))\ \sqcup\ 6(A_2(-))
\)，且两套 \(A_2(\pm)\) 仅由“两个非零分量同号/异号”的符号条件区分。
因此该分裂在字典层面不引入额外选择参数，可作为 \(\mathfrak{su}(3)\) 型子数据的内生二分机制。
\end{remark}

\paragraph{可复核接口：NAP 选择器与 \texorpdfstring{$m=11$}{m=11} 的 \texorpdfstring{$\ZZ_{34}$}{Z34} Fourier 分解证书}
小节 \ref{subsec:bdry-tower-zeck-gut} 在 Zeckendorf--GUT 签名字典之上引入了一个新的离散选择律：
要求非阿贝尔共振层 \(\{6,8\}\) 必须作为独立层保留（NAP），并在所有简单李代数中最小化 \(D=\dim\mathfrak g\)。
该选择器把 \(SO(10)\) 作为最小解离散地挑出，并给出“\(4\to 11\)”的单次签名替换与 \(34-1=33\) 的刚性计数接口。
与此同时，\(m=11\) 的 \(34\) 个边界模态在规范编号下携带 \(\ZZ_{34}\) 平移作用，从而强制一个
\(34=1\oplus 33\) 且 \(33=1\oplus 16\times 2\) 的实表示分解模板（为后续投影门/核谱的谱并群提供可审计靶标）。
\begin{itemize}
  \item \textbf{NAP 选择器脚本：}\texttt{\detokenize{scripts/exp_nap_selector_so10.py}}。
  \item \textbf{NAP LaTeX 生成物：}\texttt{\detokenize{sections/generated/tab_nap_selector_so10.tex}}，
  \texttt{\detokenize{sections/generated/eq_nap_selector_so10.tex}}。
  \item \textbf{NAP 审计 JSON：}\texttt{\detokenize{artifacts/export/nap_selector_so10.json}}（含扫描参数、候选表与最小解证书）。
  \item \textbf{\(\ZZ_{34}\) Fourier 分解脚本：}\texttt{\detokenize{scripts/exp_m11_boundary_z34_fourier_decomposition.py}}。
  \item \textbf{\(\ZZ_{34}\) LaTeX 生成物：}\texttt{\detokenize{sections/generated/eq_m11_boundary_z34_decomposition.tex}}。
  \item \textbf{\(\ZZ_{34}\) 审计 JSON：}\texttt{\detokenize{artifacts/export/m11_boundary_z34_fourier_decomposition.json}}（含规范编号 \(j\mapsto (v_7,u_{11})\) 与正交/旋转校验残差）。
\end{itemize}

\begin{theorem}[\texorpdfstring{$\ZZ_{34}$}{Z34} 平移作用的实不可约分解唯一性]\label{thm:m11-z34-real-irrep-decomposition}
\leanverified{paper\_m11\_z34\_real\_irrep\_decomposition}
设 \(V_{\RR}:=\RR^{\ZZ_{34}}\) 为 \(34\) 维实向量空间，取标准基 \(\{e_k\}_{k\in\ZZ_{34}}\)。
令 \(\ZZ_{34}\) 的生成元 \(g\) 以循环置换作用
\[
g\cdot e_k=e_{k+1}\qquad(k\in\ZZ_{34})
\]
定义 \(V_{\RR}\) 上的实置换表示。
则存在（并在 \(\ZZ_{34}\)-等变同构意义下唯一）直和分解
\begin{equation}\label{eq:m11-z34-real-irrep-sum}
\boxed{\;
V_{\RR}\ \cong\ \mathbf 1_{\mathrm{triv}}\ \oplus\ \mathbf 1_{-}\ \oplus\ \bigoplus_{k=1}^{16}W_k\;}
\end{equation}
其中 \(\mathbf 1_{\mathrm{triv}}\) 为平凡一维表示，\(\mathbf 1_{-}\) 为唯一非平凡实一维表示（满足 \(g\mapsto -1\)，对应角色指数 \(17\)），而每个 \(W_k\) 为二维不可约实表示，在其上 \(g\) 以角度 \(2\pi k/34\) 的平面旋转作用。
特别地，令增广映射 \(\varepsilon:V_{\RR}\to\RR\) 由
\(
\varepsilon(\sum_k a_ke_k)=\sum_k a_k
\)
给出，并记 \(I:=\ker\varepsilon\)。
则
\[
V_{\RR}\ \cong\ \mathbf 1_{\mathrm{triv}}\oplus I,
\qquad
\dim I=33,
\qquad
I\ \cong\ \mathbf 1_{-}\ \oplus\ \bigoplus_{k=1}^{16}W_k,
\]
从而维数与不可约型在实数域上被完全锁定为
\[
34=1\oplus 33,
\qquad
33=1\oplus 16\times 2.
\]
\end{theorem}
\begin{proof}
复化后 \(V_{\CC}:=V_{\RR}\otimes_{\RR}\CC\cong \CC[\ZZ_{34}]\) 为正则表示。
对循环群，\(\CC[\ZZ_{34}]\) 按角色分解为
\[
V_{\CC}\cong \bigoplus_{j=0}^{33}\chi_j,\qquad \chi_j(g)=\exp\!\left(\frac{2\pi i j}{34}\right).
\]
其中 \(j=0\) 给出平凡角色；由于 \(34\) 为偶数，\(j=17\) 给出实角色 \(\chi_{17}(g)=-1\)。
对每个 \(1\le j\le 16\)，\(\chi_j\) 与 \(\chi_{34-j}\) 构成共轭对；将该共轭对在 \(\RR\) 上合并，得到二维不可约实表示 \(W_j\)。
计数得出 \eqref{eq:m11-z34-real-irrep-sum}。
\par
增广映射 \(\varepsilon\) 为 \(\ZZ_{34}\)-等变映射，其像为 \(\mathbf 1_{\mathrm{triv}}\)。
由于 \(\mathbf 1_{\mathrm{triv}}\) 在 \(V_{\RR}\) 中出现重数为 \(1\)，可取 \(\varepsilon\) 的任一 \(\ZZ_{34}\)-等变右逆得到 \(\ZZ_{34}\)-等变分裂
\(
V_{\RR}\cong \mathbf 1_{\mathrm{triv}}\oplus I
\)。
由 \eqref{eq:m11-z34-real-irrep-sum} 中的角色计数可知：去掉 \(\mathbf 1_{\mathrm{triv}}\) 后恰保留 \(\mathbf 1_-\) 与全部 \(W_k\)，从而得到 \(I\) 的分解与 \(33=1\oplus 16\times 2\)。证毕。
\end{proof}

\begin{corollary}[\texorpdfstring{$32$}{32} 维旋转部分的内禀复结构]\label{cor:m11-z34-rotation-part-complex-structure}
\leanverified{paper\_m11\_z34\_rotation\_part\_complex\_structure}
在定理 \ref{thm:m11-z34-real-irrep-decomposition} 的分解中，令
\[
W:=\bigoplus_{k=1}^{16}W_k,
\qquad \dim_{\RR}W=32.
\]
则存在一个 \(\ZZ_{34}\)-不变的线性算子 \(J:W\to W\) 满足 \(J^2=-\mathrm{id}\)。
因此 \(W\) 作为复向量空间同构于 \(\CC^{16}\)，并且 \(\ZZ_{34}\) 的作用在该复结构下为复线性作用。
\end{corollary}
\begin{proof}
对每个二维不可约实表示 \(W_k\)，生成元 \(g\) 的作用在适当实基下为角度 \(2\pi k/34\) 的平面旋转。
在该基下令 \(J_k\) 为逆时针 \(90^\circ\) 旋转，则 \(J_k^2=-\mathrm{id}\)，且 \(J_k\) 与任意旋转矩阵可交换，尤其与 \(g\) 的作用可交换。
定义 \(J:=\bigoplus_{k=1}^{16}J_k\)，则 \(J\) 与 \(\ZZ_{34}\) 的作用交换并满足 \(J^2=-\mathrm{id}\)。
由此 \(W\) 获得 \(\ZZ_{34}\)-不变复结构，故 \(W\cong_{\CC}\CC^{16}\)。证毕。
\end{proof}

\begin{theorem}[window-\(11\) 的 \texorpdfstring{$\ZZ_{34}$}{Z34} 边界 uplift：内禀 \texorpdfstring{$\mathrm{SU}(16)$}{SU(16)} 作用]\label{thm:m11-z34-uplift-intrinsic-su16}
\leanverified{paper\_m11\_z34\_uplift\_intrinsic\_su16}
沿用定理 \ref{thm:m11-z34-real-irrep-decomposition} 与推论 \ref{cor:m11-z34-rotation-part-complex-structure} 的记号，令
\(
W=\bigoplus_{k=1}^{16}W_k
\)
并取任一 \(\ZZ_{34}\)-等变复结构 \(J\) 使 \(W\cong_{\CC}\CC^{16}\)。
则 \(\ZZ_{34}\) 在 \((W,J)\) 上诱导一个复线性表示
\[
\rho:\ZZ_{34}\hookrightarrow U(W,J),
\]
并且存在与 \(J\) 相容的 Hermite 内积，使得 \(\rho(\ZZ_{34})\subseteq \mathrm{SU}(W,J)\cong \mathrm{SU}(16)\)。
此外，\(\rho\) 为单射。
\end{theorem}
\begin{proof}
对每个 \(k\in\{1,\dots,16\}\)，在 \(W_k\) 上令 \(\rho_k(g)\) 为生成元作用，其为旋转角 \(\theta_k=2\pi k/34\in(0,\pi)\) 的平面旋转。
取与 \(\rho_k(g)\) 相容的标准内积，并定义
\[
J_k:=\frac{1}{\sin\theta_k}\Bigl(\rho_k(g)-\cos\theta_k\,\mathrm{Id}\Bigr)\in \mathrm{End}_{\RR}(W_k).
\]
则 \(\rho_k(g)=\cos\theta_k\,\mathrm{Id}+\sin\theta_k\,J_k\)，从而 \(J_k^2=-\mathrm{Id}\)；并且 \(J_k\) 属于由 \(\rho_k(g)\) 生成的实交换子代数，故与 \(\rho_k(\ZZ_{34})\) 对易。
令 \(J:=\bigoplus_{k=1}^{16}J_k\)，得到 \(\ZZ_{34}\)-等变复结构。
\par
在复结构 \(J\) 下，\(g\) 在第 \(k\) 个复一维分量上按标量 \(e^{\mathrm{i}\theta_k}=e^{2\pi \mathrm{i}k/34}\) 作用。
因此
\[
\det_{\CC}\rho(g)=\prod_{k=1}^{16}e^{2\pi \mathrm{i}k/34}
=\exp\!\left(\frac{2\pi \mathrm{i}}{34}\sum_{k=1}^{16}k\right)
=\exp\!\left(\frac{2\pi \mathrm{i}}{34}\cdot 136\right)=1,
\]
从而 \(\rho(\ZZ_{34})\subseteq \mathrm{SU}(W,J)\)。
若 \(\rho(g^t)=\mathrm{Id}\)，则取 \(k=1\) 得 \(e^{2\pi \mathrm{i}t/34}=1\)，从而 \(34\mid t\)，故 \(\rho\) 单射。证毕。
\end{proof}

\begin{proposition}[\texorpdfstring{$\ZZ_{34}$}{Z34} 相位自由度的交换子刚性与 \texorpdfstring{$U(1)^{16}$}{U(1)^16} 的内禀出现]\label{prop:m11-z34-commutant-u1-16}
沿用定理 \ref{thm:m11-z34-uplift-intrinsic-su16} 的设定。
则 \(G=\ZZ_{34}\) 在 \(W=\bigoplus_{k=1}^{16}W_k\) 上的实交换子代数满足实代数同构
\[
\mathrm{End}_{\RR,G}(W)\ :=\ \{T\in \mathrm{End}_{\RR}(W):\ T\rho(h)=\rho(h)T,\ \forall h\in G\}
\ \cong\ \CC^{16}.
\]
若进一步取与 \(J\) 相容的 Hermite 内积，则中心化子满足
\[
Z_{U(W,J)}(\rho(G))\ \cong\ U(1)^{16}.
\]
\leanverified{paper\_m11\_z34\_commutant\_u1\_16}
\end{proposition}
\begin{proof}
由 \(W=\bigoplus_{k=1}^{16}W_k\) 且各 \(W_k\) 两两非同构的二维不可约实 \(G\)-模（旋转角不同），有
\[
\mathrm{End}_{\RR,G}(W)\ \cong\ \prod_{k=1}^{16}\mathrm{End}_{\RR,G}(W_k).
\]
对每个 \(k\)，\(W_k\) 为二维旋转不可约表示，其交换子由 \(\mathrm{Id}\) 与复结构 \(J_k\) 生成，故 \(\mathrm{End}_{\RR,G}(W_k)\cong \RR[J_k]\cong \CC\)（作为实代数）。
乘积即得 \(\CC^{16}\)。
\par
在给定 \(J\)-相容 Hermite 内积下，每个 \(\CC\) 因子中保范元对应单位圆 \(U(1)\)，从而中心化子为 \(U(1)^{16}\)。证毕。
\end{proof}

\begin{theorem}[偶边界与单层替换约束下的极小简单解：\texorpdfstring{$m=11$}{m=11} 与 \texorpdfstring{$\mathfrak{so}(10)$}{so(10)}]\label{thm:bdry-even-single-layer-minimal-so10}
\leanverified{paper\_bdry\_even\_single\_layer\_minimal\_so10}
令 \(|X_m^{\mathrm{bdry}}|=F_{m-2}\) 为边界层计数，并固定两层 \(\{6,8\}\)。
对 \(m>8\) 令
\(
\Sigma(m):=\{6,8,m\}
\)
并令签名维数
\(
D(\Sigma(m)):=F_{m-2}+F_{6}+F_{4}=F_{m-2}+11
\)。
若要求：
\begin{enumerate}
  \item \(F_{m-2}\) 为偶数；
  \item 存在紧简单李代数 \(\mathfrak g\) 满足 \(\dim\mathfrak g=D(\Sigma(m))\)，且在满足 (1) 的所有 \(m>8\) 中该维数取最小；
\end{enumerate}
则唯一可行解为 \(m=11\)，且 \(\mathfrak g\cong \mathfrak{so}(10)\)（\(\dim=45\)）。
\end{theorem}
\begin{proof}
由 Fibonacci 偶奇性，\(F_n\) 为偶当且仅当 \(3\mid n\)，故 (1) 等价于 \(3\mid (m-2)\)。
在 \(m>8\) 中满足该条件的最小者为 \(m=11\)，并且 \(F_{9}=34\)，从而
\(
D(\Sigma(11))=34+11=45
\)。
在紧简单李代数分类中，维数 \(45\) 的简单代数同构型唯一，为 \(\mathfrak{so}(10)\)。
证毕。
\end{proof}

\begin{corollary}[家族数锁定分支上的 \(16\) 个二维旋转平面]\label{cor:m11-z34-sixteen-rotation-planes-from-family-lock}
在定理 \ref{thm:terminal-family-uplift-lock} 所锁定的 \(F_9=34\) 分支上，由定理 \ref{thm:bdry-delta34-m11-uniqueness} 得唯一分辨率 \(m=11\) 满足 \(|X_{11}^{\mathrm{bdry}}|=34\)。
若在该边界层上取 \(\ZZ_{34}\) 的平移作用作为可观测对称接口，则 \(\RR^{X_{11}^{\mathrm{bdry}}}\) 的实不可约分解中非平凡二维旋转块的个数恰为 \(16\)。
\end{corollary}
\begin{proof}
由定理 \ref{thm:bdry-delta34-m11-uniqueness} 得 \(|X_{11}^{\mathrm{bdry}}|=34\)，从而可将边界平移对称表述为 \(\ZZ_{34}\) 的正则置换作用。
应用定理 \ref{thm:m11-z34-real-irrep-decomposition} 即得二维旋转块数为 \((34-2)/2=16\)。证毕。
\end{proof}

\begin{proposition}[\texorpdfstring{$\QQ[\ZZ_{34}]$}{Q[Z34]} 的分解与两份 \texorpdfstring{$17$}{17} 次割圆域]\label{prop:m11-qz34-cyclotomic-decomposition}
\leanverified{paper\_m11\_qz34\_cyclotomic\_decomposition}
令
\[
A:=\QQ[\ZZ_{34}]\cong \QQ[x]/(x^{34}-1).
\]
则存在有理代数同构
\begin{equation}\label{eq:m11-qz34-cyclotomic-decomposition}
\boxed{\;
\QQ[\ZZ_{34}]
\ \cong\
\QQ\ \oplus\ \QQ\ \oplus\ K\ \oplus\ K,
\qquad
K:=\QQ(\zeta_{17})\;}
\end{equation}
其中 \(\zeta_{17}\) 为原始 \(17\) 次单位根，且 \([K:\QQ]=\deg\Phi_{17}=16\)。
\end{proposition}
\begin{proof}
由割圆分解 \(x^{34}-1=\prod_{d\mid 34}\Phi_d(x)\)，并注意 \(34\) 的正因子为 \(1,2,17,34\)，得
\[
x^{34}-1=\Phi_1(x)\Phi_2(x)\Phi_{17}(x)\Phi_{34}(x).
\]
在 \(\QQ[x]\) 上这些因子两两互素，故由中国剩余定理
\[
\QQ[x]/(x^{34}-1)\cong \bigoplus_{d\mid 34}\QQ[x]/(\Phi_d(x))
\cong \QQ\oplus \QQ\oplus \QQ(\zeta_{17})\oplus \QQ(\zeta_{34}).
\]
又因 \(\zeta_{34}=-\zeta_{17}\) 且 \(-1\in\QQ\)，故 \(\QQ(\zeta_{34})=\QQ(\zeta_{17})=:K\)。
度数 \([K:\QQ]=\deg\Phi_{17}=16\) 为标准事实。证毕。
\end{proof}

\begin{corollary}[内禀 \texorpdfstring{$C_{16}$}{C16} 自同构：\(16\) 的有理层强制显影]\label{cor:m11-qz34-galois-c16}
\leanverified{paper\_m11\_qz34\_galois\_c16}
在命题 \ref{prop:m11-qz34-cyclotomic-decomposition} 的记号下，
\[
\mathrm{Gal}(K/\QQ)\cong (\ZZ/17\ZZ)^\times\cong C_{16}.
\]
因此在分解 \(\QQ[\ZZ_{34}]\cong \QQ^2\oplus K^2\) 中，每个 \(K\) 分量均携带一个内禀的 \(C_{16}\) 作用；该作用与 \(\QQ^2\) 分量正交分离，并在 \(16\) 维有理不可约块上给出唯一的循环生成方式。
\end{corollary}
\begin{proof}
对素数 \(p\)，有 \(\mathrm{Gal}(\QQ(\zeta_p)/\QQ)\cong (\ZZ/p\ZZ)^\times\)。
当 \(p=17\) 时，\((\ZZ/17\ZZ)^\times\) 为有限域乘法群，必为循环群，阶为 \(16\)。
将自同构 \(\sigma_a(\zeta_{17})=\zeta_{17}^a\) 逐分量作用即可。证毕。
\end{proof}

\begin{proposition}[增广理想的等变唯一性：余维 \(1\) 的平凡商强制锁定 \(33\) 维子模]\label{prop:m11-z34-augmentation-ideal-unique}
\leanverified{paper\_m11\_z34\_augmentation\_ideal\_unique}
沿用定理 \ref{thm:m11-z34-real-irrep-decomposition} 的记号。
设 \(J\subset V_{\RR}\) 为 \(\ZZ_{34}\)-子模，满足商模 \(V_{\RR}/J\) 同构于平凡表示 \(\mathbf 1_{\mathrm{triv}}\)。
则必有 \(J=I=\ker\varepsilon\)。
因此满足“余维 \(1\) 且商为平凡表示”的 \(33\) 维 \(\ZZ_{34}\)-子模在等变意义下无自由度。
\end{proposition}
\begin{proof}
由假设存在 \(\ZZ_{34}\)-等变满射 \(\ell:V_{\RR}\twoheadrightarrow \RR\) 使 \(J=\ker\ell\)。
而
\(
\mathrm{Hom}_{\ZZ_{34}}(V_{\RR},\mathbf 1_{\mathrm{triv}})
\cong
(V_{\RR})^{\ZZ_{34}}
\)
与不变子空间对偶同构，维数为 \(1\)（由正则表示中平凡因子重数为 \(1\)）。
故 \(\ell\) 与 \(\varepsilon\) 成比例，从而 \(\ker\ell=\ker\varepsilon=I\)。证毕。
\end{proof}

\begin{corollary}[维数补余的表示论锁定：\texorpdfstring{$33=45-12$}{33=45-12} 的唯一模实现]\label{cor:m11-z34-unique-33-module}
\leanverified{paper\_m11\_z34\_unique\_33\_module}
若在 \(\ZZ_{34}\) 平移对称下要求一个 \(33\) 维结构与 \(V_{\RR}\) 的平凡商方向正交（即商模为 \(\mathbf 1_{\mathrm{triv}}\)），
则该 \(33\) 维结构必与增广理想 \(I\) 同构，从而自动继承 \(33=1\oplus 16\times 2\) 的实不可约型。
\end{corollary}
\begin{proof}
由命题 \ref{prop:m11-z34-augmentation-ideal-unique}，满足条件的余维 \(1\) 子模唯一，即为 \(I\)。
再由定理 \ref{thm:m11-z34-real-irrep-decomposition} 得到 \(I\) 的不可约型。证毕。
\end{proof}

\begin{corollary}[window-\(6\) 局域 uplift 与 \texorpdfstring{$\ZZ_{34}$}{Z34} 平移证书的割圆强迫]\label{cor:window6-local-uplift-forces-cyclotomic17}
\leanverified{paper\_window6\_local\_uplift\_forces\_cyclotomic17}
在 window-\(6\) 的“局域几何 uplift”协议口径下，推论 \ref{cor:terminal-window6-local-uplift-admissibility} 给出除平凡位移外仅允许 \(\delta=\pm 34\)。
若进一步在 \(m=11\) 边界模态上存在 \(\ZZ_{34}\) 平移作用（见 \S\ref{subsec:group-unification-audit-pointers} 的审计接口），
则有理层上的不可约分量必包含 \(K=\QQ(\zeta_{17})\)（度数 \(16\)），并伴随内禀 \(C_{16}\) 伽罗瓦对称（推论 \ref{cor:m11-qz34-galois-c16}）。
\end{corollary}
\begin{proof}
局域 uplift 的唯一整数指纹把位移尺度锁定为 \(34\)。
而 \(\ZZ_{34}\) 平移作用等价于引入群代数 \(\QQ[\ZZ_{34}]\) 的表示接口；由命题 \ref{prop:m11-qz34-cyclotomic-decomposition}，
\(\QQ[\ZZ_{34}]\) 的有理不可约分量中必出现两份同构的割圆域 \(K=\QQ(\zeta_{17})\)，其伽罗瓦群为 \(C_{16}\)（推论 \ref{cor:m11-qz34-galois-c16}）。
证毕。
\end{proof}

\paragraph{\texorpdfstring{$\delta$}{δ}-共振闭包、\texorpdfstring{$m=11$}{m=11} 电荷层与 \texorpdfstring{$\mathrm{Spin}(10)$}{Spin(10)} 接口的进一步刚性}

\begin{theorem}[\texorpdfstring{$\delta$}{δ}-共振闭包下的唯一极小简单统一型]\label{thm:m11-delta-closure-unique-minimal-so10}
\leanverified{paper\_m11\_delta\_closure\_unique\_minimal\_so10}
在边界塔计数律
\[
|X_m|=F_{m+2},\qquad |X_m^{\mathrm{bdry}}|=F_{m-2}
\]
下，设 \(\mathfrak g\) 为简单李代数，且其分辨率签名 \(\Sigma(\mathfrak g)\) 满足：
\begin{enumerate}
  \item \(\{6,8\}\subseteq \Sigma(\mathfrak g)\)（保留两条非阿贝尔共振层）；
  \item 存在 \(m_\delta\in \Sigma(\mathfrak g)\) 使 \(|X_{m_\delta}^{\mathrm{bdry}}|=\delta\)，其中 \(\delta=34\) 为 window-\(6\) 的 boundary 二层 uplift 非零差值。
\end{enumerate}
则 \(m_\delta=11\)，从而
\[
\{6,8,11\}\subseteq \Sigma(\mathfrak g),
\qquad
\dim\mathfrak g\ge |X_6^{\mathrm{bdry}}|+|X_8^{\mathrm{bdry}}|+|X_{11}^{\mathrm{bdry}}|
=3+8+34=45.
\]
并且等号成立当且仅当 \(\Sigma(\mathfrak g)=\{6,8,11\}\)。
若再要求在满足上述两条的简单李代数中维数取极小，则唯一可能为
\[
\mathfrak g\cong \mathfrak{so}(10).
\]
\end{theorem}
\begin{proof}
由定理 \ref{thm:bdry-delta34-m11-uniqueness}，条件 \(|X_m^{\mathrm{bdry}}|=\delta=34\) 的解唯一且 \(m=11\)。
故 \(\Sigma(\mathfrak g)\) 至少含 \(\{6,8,11\}\)。
在本文签名口径下，
\[
\dim\mathfrak g=\sum_{m\in\Sigma(\mathfrak g)}|X_m^{\mathrm{bdry}}|
\]
且各项非负，因而立得维数下界 \(45\)。
若等号成立，则不能再含其他层，故 \(\Sigma(\mathfrak g)=\{6,8,11\}\)；反之该签名即给出 \(45\)。
\par
下面分类 \(45\) 维简单型。
若 \(\mathfrak g\cong\mathfrak{su}(n)\)，则 \(n^2-1=45\) 无整数解；
若 \(\mathfrak g\cong\mathfrak{sp}(n)\)，则 \(n(2n+1)=45\) 无整数解；
若 \(\mathfrak g\cong\mathfrak{so}(n)\)，则 \(\frac{n(n-1)}{2}=45\) 给出 \(n=10\)；
例外型维数不取 \(45\)。
故唯一可能为 \(\mathfrak{so}(10)\)。证毕。
\end{proof}

\begin{corollary}[内禀手征读出强制色荷层保留]\label{cor:m11-intrinsic-chirality-forces-color-layer}
若采用协议内禀二值手征准则（定理 \ref{thm:window6-chirality-anchor-minimal}），则手征锚点被唯一固定在 \(m=6\)；\(m=2,4\) 仅提供维数接口而不产生内禀二值手征通道。
在此基础上，一旦要求实现
\[
\delta=|X_8|-|X_6|=34
\]
的 uplift 差值（定理 \ref{thm:bdry-delta34-m11-uniqueness}），则 \(m=8\) 层不可删除。
由 \(|X_8^{\mathrm{bdry}}|=F_6=8\) 及边界塔最小共振三元组对齐（推论 \ref{cor:bdry-tower-minimal-resonance-triple-sm}），可见 \(\mathfrak{su}(3)\) 色荷层在该框架中不可避免。
\end{corollary}
\begin{proof}
第一部分即定理 \ref{thm:window6-chirality-anchor-minimal}。
若删除 \(m=8\) 层，则 \(|X_8|\) 不再参与同一分辨率差分口径，\(|X_8|-|X_6|=34\) 的结构恒等式无法保持。
另一方面，\(m=8\) 的边界计数固定为 \(8\)，并在最小共振三元组中对应 \(\mathfrak{su}(3)\) 的维数槽位。
故内禀手征与 \(\delta\)-差值并置时，\(\mathfrak{su}(3)\) 层为必需。证毕。
\end{proof}
\leanverified{paper\_m11\_intrinsic\_chirality\_forces\_color\_layer}

\begin{proposition}[\texorpdfstring{$m=11$}{m=11} 电荷层的规范 \texorpdfstring{$16$}{16} 维复模与单位行列式]\label{prop:m11-canonical-16-complex-module-det-one}
\leanverified{paper\_m11\_canonical\_16\_complex\_module\_det\_one}
在 \(X_{11}^{\mathrm{bdry}}\cong \ZZ_{34}\) 的规范编号下，令 \(T:e_j\mapsto e_{j+1}\) 为平移算子，并沿用 \eqref{eq:m11-z34-real-irrep-sum} 的分解。
记
\[
W:=\bigoplus_{k=1}^{16}W_k.
\]
则 \(W\) 具有规范复结构，\(\dim_{\CC}W=16\)，且 \(T\) 在 \(W\) 上诱导复线性自同构。
若取复基 \(v_k=C_k+\mathrm{i}S_k\)（\(k=1,\dots,16\)），并设 \(\zeta:=\exp(2\pi \mathrm{i}/34)\)，则
\[
T(v_k)=\zeta^k v_k,\qquad
\det_{\CC}(T|_W)=\prod_{k=1}^{16}\zeta^k=\zeta^{136}=1.
\]
故 \(T|_W\in \mathrm{SL}(W)\)。
\end{proposition}
\begin{proof}
由定理 \ref{thm:m11-z34-real-irrep-decomposition} 与推论 \ref{cor:m11-z34-rotation-part-complex-structure}，
\(W=\bigoplus_{k=1}^{16}W_k\) 具有 \(\ZZ_{34}\)-等变复结构，且 \(\dim_{\CC}W=16\)。
在每个二维旋转平面 \(W_k\) 上，\(T\) 的特征值为 \(e^{2\pi \mathrm{i}k/34}=\zeta^k\)，于是
\[
\det_{\CC}(T|_W)=\prod_{k=1}^{16}\zeta^k
=\zeta^{\sum_{k=1}^{16}k}
=\zeta^{136}
=\zeta^{4\cdot 34}=1.
\]
因此 \(T|_W\) 落在 \(\mathrm{SL}(W)\)。证毕。
\end{proof}

\begin{theorem}[\texorpdfstring{$m=11$}{m=11} 平移的二面体提升与 Dirac 型双重化]\label{thm:m11-dihedral-lift-dirac-doubling}
\leanverified{paper\_m11\_dihedral\_lift\_dirac\_doubling}
在 \(\RR^{\ZZ_{34}}\) 上定义
\[
T(e_j):=e_{j+1},\qquad S(e_j):=e_{-j}\qquad (j\in\ZZ_{34}).
\]
则
\[
STS=T^{-1},
\]
从而 \(\langle T,S\rangle\cong D_{34}\)（阶 \(68\) 的二面体群）。
限制到命题 \ref{prop:m11-canonical-16-complex-module-det-one} 的 \(W\) 上，有
\[
S(C_k)=C_k,\qquad S(S_k)=-S_k,\qquad S(v_k)=\overline{v_k}.
\]
于是作为复表示成立同构
\[
W_{\RR}\otimes_{\RR}\CC\ \cong\ W\oplus \overline{W},
\]
即该 \(16\) 维复模在实化后出现规范的共轭双重化。
\end{theorem}
\begin{proof}
先验证群关系：
\[
STS(e_j)=ST(e_{-j})=S(e_{-j+1})=e_{j-1}=T^{-1}(e_j),
\]
并且 \(S^2=\mathrm{Id}\)、\(T^{34}=\mathrm{Id}\)，故生成关系即二面体关系，得到 \(\langle T,S\rangle\cong D_{34}\)。
\par
再看 Fourier 平面基。
由
\[
C_k=\sqrt{\frac{2}{34}}\sum_{j=0}^{33}\cos\!\left(\frac{2\pi kj}{34}\right)e_j,\qquad
S_k=\sqrt{\frac{2}{34}}\sum_{j=0}^{33}\sin\!\left(\frac{2\pi kj}{34}\right)e_j,
\]
代入 \(j\mapsto -j\) 可得 \(S(C_k)=C_k\)、\(S(S_k)=-S_k\)，从而
\[
S(v_k)=S(C_k+\mathrm{i}S_k)=C_k-\mathrm{i}S_k=\overline{v_k}.
\]
因此 \(S\) 在 \(W\) 上对应复共轭，标准实化分解给出
\(
W_{\RR}\otimes_{\RR}\CC\cong W\oplus\overline{W}
\)。
证毕。
\end{proof}

\begin{corollary}[行列式线丛平凡性与一阶离散异常消失]\label{cor:m11-det-line-trivial-c1-zero}
\leanverified{paper\_m11\_det\_line\_trivial\_c1\_zero}
令 \(E\to B\ZZ_{34}\) 为由命题 \ref{prop:m11-canonical-16-complex-module-det-one} 的复表示 \(W\) 诱导的平坦复向量丛。
则其行列式线丛 \(\det E\) 平凡，且
\[
c_1(E)=c_1(\det E)=0\in H^2(B\ZZ_{34};\ZZ)\cong \ZZ_{34}.
\]
\end{corollary}
\begin{proof}
行列式线丛由行列式特征给出，而命题 \ref{prop:m11-canonical-16-complex-module-det-one} 已证
\(
\det_{\CC}(T|_W)=1
\)。
故 \(\det E\) 为平凡线丛，进而 \(c_1(\det E)=0\)。
由 \(c_1(E)=c_1(\det E)\) 得结论。证毕。
\end{proof}

\begin{conjecture}[\texorpdfstring{$m=11$}{m=11} 的 \texorpdfstring{$16$}{16} 模作为 \texorpdfstring{$\mathrm{Spin}(10)$}{Spin(10)} 半自旋离散骨架]\label{conj:m11-spin10-halfspin-discrete-skeleton}
在定理 \ref{thm:m11-delta-closure-unique-minimal-so10} 与命题 \ref{prop:m11-canonical-16-complex-module-det-one} 的框架下，设
\[
\iota:\ \ZZ_{34}\rtimes \ZZ_2\ \hookrightarrow\ \mathrm{Spin}(10)
\]
为单射同态，其中 \(\ZZ_2\) 对 \(\ZZ_{34}\) 的作用取反演 \(r\mapsto -r\)（对应定理 \ref{thm:m11-dihedral-lift-dirac-doubling} 的 \(D_{34}\) 关系）。
猜想存在某一半自旋表示 \(S^+\) 使：
\begin{enumerate}
  \item \(S^+|_{\iota(\ZZ_{34})}\cong W\)；
  \item 在
  \[
  \mathfrak g_{\mathrm{SM}}\cong
  \mathfrak u(1)\oplus\mathfrak{su}(2)\oplus\mathfrak{su}(3)
  \]
  的限制下，\(W\) 的分解与一代标准模型 Weyl 费米子表示内容同构（包括色重数与弱同位旋重数）；
  \item 超荷由 \(\ZZ_{34}\) 特征指数实现离散量子化，并与 (2) 的分解兼容。
\end{enumerate}
\end{conjecture}

\endinput


\begin{conjecture}[boundary uplift 位移谱与 reset-event 归返间隙谱的同一性]\label{conj:bdry-uplift-delta-equals-reset-gaps}
对 \(m\ge 6\)，令 \(\Delta_m^{\mathrm{bdry}}\subset\ZZ\) 表示 boundary 扇区在 dyadic uplift 下的位移集合（例如表 \ref{tab:foldbin_boundary_lift_m6_m8} 给出 \(m\in\{6,7,8\}\) 的证书输出）。
令 \(\mathcal R_m\) 为复位事件集合（定理 \ref{thm:terminal-reset-events-sturmian}），并记其归返间隙谱为 \(\mathrm{Gap}(\mathcal R_m)\)（其中已知闭式为 \(\{F_{m+3},F_{m+4}\}\)）。
则猜想：对一切 \(m\ge 7\)，成立集合恒等式
\[
\Delta_m^{\mathrm{bdry}}\setminus\{0\}=\mathrm{Gap}(\mathcal R_m)=\{F_{m+3},F_{m+4}\},
\]
且 \(m=6\) 为唯一退化例外：\(\Delta_6^{\mathrm{bdry}}\setminus\{0\}=\{F_9\}\subsetneq \{F_9,F_{10}\}\)。
\end{conjecture}

\begin{conjecture}[偶边界窗口的增广理想与统一候选的函子化匹配公理]\label{conj:profinite-gut-augmentation-ideal-functor}
设 \(m\equiv 2\pmod 3\)，从而 \(|X_m^{\mathrm{bdry}}|=F_{m-2}\) 为偶数。
令 \(N:=|X_m^{\mathrm{bdry}}|\) 并记 \(G_m:=\ZZ_N\)。
将 \(\RR[G_m]\) 视为正则表示，并记其增广理想
\[
I_m:=\ker\!\bigl(\varepsilon:\RR[G_m]\to\RR\bigr).
\]
预期存在一个从“偶边界窗口”数据到紧致李群表示范畴的函子 \(\mathcal F\)，使得：
\begin{enumerate}
  \item \(\mathcal F(m)\) 的最小非平凡表示包含一个由 \(I_m\) 的 Fourier 分解型确定的相位骨架，其交换子代数分解为若干 \(\CC\) 因子直积；
  \item 在 \(m=11\) 处，\(\mathcal F(11)\) 的最小简单部分与 \(\mathrm{Spin}(10)\)（等价地 \(\mathfrak{so}(10)\)）兼容，并且 \(I_{11}\) 的分解型与半自旋分支之间满足由 \(\mathbf 1_{-}\) 刚性因子锁定的相容条件；
  \item 对一般 \(m\equiv 2\pmod 3\)，该相容条件诱导一条可计算的离散筛法，用以从增广理想的分解型算术不变量中筛选统一候选。
\end{enumerate}
\end{conjecture}

\begin{theorem}[\texorpdfstring{$\ZZ_{34}$}{Z34} 的 \texorpdfstring{$\mathbb F_{17}$}{F17} 商影上的和积逃逸与偶旋转块证书]\label{distill:bourgain-sum_product_obstruction_classification-z34-f17-quotient-expansion-even-rotation-escape}
令
\[
\pi_{17}:\ZZ_{34}\longrightarrow \mathbb F_{17}
\]
为模 \(17\) 商映射。取非空集合 \(B\subset\ZZ_{34}\)，并记
\[
C:=\pi_{17}(B)\subset\mathbb F_{17}.
\]
设 \(2\le |C|\le16\)。固定素域 \(\mathbb F_{17}\) 上的一组 Bourgain--Katz--Tao 和积常数 \(c_{17}>0\)、\(\varepsilon_{17}>0\)，使得对任意 \(E\subset\mathbb F_{17}\) 且 \(2\le |E|\le16\) 有
\[
\max\{|E+E|,|E\cdot E|\}\ge c_{17}|E|^{1+\varepsilon_{17}}.
\]
若某个 \(K\ge1\) 同时满足
\[
|\pi_{17}(B+B)|\le K|C|,
\qquad
|\pi_{17}(B\cdot B)|\le K|C|,
\]
则必有
\[
K\ge c_{17}|C|^{\varepsilon_{17}}.
\]
等价地，若 \(K<c_{17}|C|^{\varepsilon_{17}}\)，则 \(\ZZ_{34}\) 窗口在 \(\mathbb F_{17}\) 商影上必发生加法逃逸或乘法逃逸：
\[
|C+C|>K|C|
\qquad\text{或}\qquad
|C\cdot C|>K|C|.
\]
进一步，令 \(E\) 为发生逃逸的一侧，即 \(E=C+C\) 或 \(E=C\cdot C\)，并定义其 \(17\)-商提升
\[
\widetilde E:=\pi_{17}^{-1}(E)\subset\ZZ_{34},
\qquad
f_E:=\mathbf 1_{\widetilde E}-\frac{|E|}{17}\mathbf 1_{\ZZ_{34}}\in V_{\RR}.
\]
则 \(f_E\) 属于定理 \ref{thm:m11-z34-real-irrep-decomposition} 的偶旋转部分
\[
f_E\in \bigoplus_{\ell=1}^{8}W_{2\ell}.
\]
并且 \(f_E=0\) 当且仅当 \(E=\mathbb F_{17}\)。因此，非饱和逃逸必在偶数二维旋转块中留下非零实 Fourier 证书；饱和逃逸则把 \(\mathbb F_{17}\) 商影推到全商环。特别地，任何只使用平凡块、符号块 \(\mathbf 1_-\) 或奇数旋转块的 group-window coincidence 证书，都不能解释这种 \(\mathbb F_{17}\) 商影上的同时低和积失败。
\end{theorem}
\begin{proof}
由于 \(\pi_{17}\) 是环同态，直接有
\[
\pi_{17}(B+B)=C+C,
\qquad
\pi_{17}(B\cdot B)=C\cdot C.
\]
故两个假设不等式等价于
\[
|C+C|\le K|C|,
\qquad
|C\cdot C|\le K|C|.
\]
对 \(C\subset\mathbb F_{17}\) 应用素域和积定理，得到
\[
c_{17}|C|^{1+\varepsilon_{17}}
\le
\max\{|C+C|,|C\cdot C|\}
\le K|C|,
\]
从而 \(K\ge c_{17}|C|^{\varepsilon_{17}}\)。这证明了和积逃逸二择一。

现在证明表示论证书。函数 \(\mathbf 1_{\widetilde E}\) 对平移 \(x\mapsto x+17\) 不变，因此其复 Fourier 展开只含在子群 \(17\ZZ_{34}\) 上平凡的角色。对 \(\ZZ_{34}\) 的角色 \(\chi_j(g)=\exp(2\pi i j/34)\)，这一条件为
\[
\chi_j(17)=\exp(\pi i j)=1,
\]
即 \(j\) 为偶数。减去平均值只去掉 \(j=0\) 的平凡分量，故剩余实向量 \(f_E\) 只落在偶数共轭角色对对应的实二维块中，也就是
\[
f_E\in\bigoplus_{\ell=1}^{8}W_{2\ell}.
\]
若 \(E=\mathbb F_{17}\)，则 \(\widetilde E=\ZZ_{34}\)，显然 \(f_E=0\)。反过来，若 \(f_E=0\)，则 \(\mathbf 1_{\widetilde E}\) 为常值函数。由于 \(\widetilde E\) 是非空集合的指示函数，只能有 \(\widetilde E=\ZZ_{34}\)，即 \(E=\mathbb F_{17}\)。

因此，非饱和的商影逃逸必然被偶旋转块检测；若逃逸饱和，则商影已经达到全 \(\mathbb F_{17}\) 层。平凡块只记录平均值，\(\mathbf 1_-\) 只记录模 \(2\) 符号方向，奇数旋转块不在 \(\pi_{17}\)-提升函数的 Fourier 支撑中，故这些块不能承载上述 \(\mathbb F_{17}\) 和积逃逸。证毕。
\end{proof}


\endinput


\begin{remark}[局域 uplift 的唯一性与 \texorpdfstring{$SU(5)$}{SU(5)}/\texorpdfstring{$E_6$}{E6} 分支的非局域性]\label{rem:window6-local-uplift-excludes-su5-e6}
由推论 \ref{cor:terminal-window6-local-uplift-admissibility}，
在保持稳定类型标签 \(F_6\) 且由 \(\mathrm{Aut}(Q_6)\) 实现的局域 uplift 中，除 \(0\) 外唯一可实现的整数位移为 \(\pm 34\)；
从而在三分支 \(\{F_8,F_9,F_{10}\}=\{21,34,55\}\) 中仅 \(F_9\) 与该局域性相容。
另一方面，定理 \ref{thm:terminal-window6-tail-three-branch} 给出 \((SU(5),SO(10),E_6)\) 的顶层 Zeckendorf 项分别对齐 \((F_8,F_9,F_{10})\)。
因此在额外施加“uplift 必须由局域几何稳定子实现”的协议约束时，
\(SU(5)\) 与 \(E_6\) 的顶层分支只能通过非局域机制出现，而 \(SO(10)\) 分支是唯一保留局域可实现性的候选。
\end{remark}

\begin{corollary}[“为何是标准模型而非循环引用”的三重刚性闭包（window-\(6\)）]\label{cor:window6-threefold-rigidity-closure-sm-so10}
在 window-\(6\) 锚点上，同时并置以下三条可审计输入：
\begin{enumerate}
  \item \textbf{几何刚性：}局域 uplift 除 \(0\) 外仅允许 \(\delta=34\)（推论 \ref{cor:terminal-window6-local-uplift-admissibility}）。
  \item \textbf{边界塔可观测性：}最小三层边界塔给出 \((1,3,8)\) 维数接口（\texttt{\detokenize{sections/generated/eq_boundary_tower_sm12.tex}}），并刚性选出 \(\mathfrak u(1)\oplus\mathfrak{su}(2)\oplus\mathfrak{su}(3)\)（推论 \ref{cor:bdry-tower-minimal-resonance-triple-sm}）。
  \item \textbf{组合刚性：}\(\mathrm{Aut}(\Gamma_6)=\{1\}\) 排除非平凡重标号（命题 \ref{prop:terminal-gamma6-rigidity} 与推论 \ref{cor:terminal-gamma6-multiplicity-rigidity}）。
\end{enumerate}
则低分辨率可见规范代数在同构意义下被规范定位为
\[
\mathfrak g_{\mathrm{vis}}\ \cong\ \mathfrak u(1)\ \oplus\ \mathfrak{su}(2)\ \oplus\ \mathfrak{su}(3),
\]
且不依赖任何先验“手工指定规范群”的输入。
进一步地，在施加 NAP 约束 \(\{6,8\}\subseteq \Sigma(\mathfrak g)\) 并最小化统一代价时，
脚本化扫描把最小简单补全离散地锁定为 \(SO(10)\)，并给出单层替换证书
\(
\Sigma_{SO(10)}=\Sigma_{\mathrm{SM}}\setminus\{4\}\cup\{11\}
\)
与 \(34-1=45-12\)（见 \eqref{eq_nap_selector_so10}）；
该替换的顶层项 \(F_9=34\) 与局域 uplift 的唯一整数指纹一致，从而在同一整数证书上闭合“几何局域性—边界塔最小性—统一补全”的非循环推理链。
\leanverified{paper\_window6\_threefold\_rigidity\_closure\_sm\_so10}
\end{corollary}
\begin{proof}
第一部分由推论 \ref{cor:terminal-window6-local-uplift-admissibility}、推论 \ref{cor:bdry-tower-minimal-resonance-triple-sm} 与命题 \ref{prop:terminal-gamma6-rigidity}（及其推论）并置即得。
第二部分由 NAP 选择器的生成等式 \eqref{eq_nap_selector_so10} 给出。证毕。
\end{proof}

\paragraph{可复核检验：三通道势—响应辛相空间的 \texorpdfstring{$\mathrm{Sp}(6)$}{Sp(6)} 各向同性点扫描}
注 \ref{rem:sp6-minimal-continuous-envelope} 给出一个把 \(B_3/C_3\) 的离散统一种子“提升为连续包络”的最小口径：在压力 \(P(\theta)\) 的 Legendre 对偶坐标下，
\((\theta;J=\nabla P)\) 诱导标准辛形式，保持三通道共轭结构的最小连续群为 \(\mathrm{Sp}(6)\)。
对应的一个严格可核对任务是：在某个“统一点”上，二次涨落尺度应趋于单尺度（各向同性），其完全可审计的分数为
\(
\mathcal{I}(\theta)=\log(\lambda_{\max}/\lambda_{\min})
\)
（见附录注 \ref{rem:sync-kernel-3d-sp6-isotropy-scan}）。
该扫描已被脚本化并纳入一键流水线：
\begin{itemize}
  \item \textbf{生成脚本：}\texttt{\detokenize{scripts/exp_sync_kernel_weighted_sp6_isotropy_scan.py}}。
  \item \textbf{LaTeX 生成物：}\texttt{\detokenize{sections/generated/tab_sync_kernel_weighted_sp6_isotropy_scan.tex}}（表 \ref{tab:sync_kernel_weighted_sp6_isotropy_scan}）。
  \item \textbf{审计 JSON：}\texttt{\detokenize{artifacts/export/sync_kernel_weighted_sp6_isotropy_scan.json}}（含扫描网格、\(\Sigma=\nabla^2P\)、特征值与最优候选点列表）。
  \item \textbf{根型分裂比率脚本：}\texttt{\detokenize{scripts/exp_sp6_isotropy_root_split_ratio.py}}（读取上条 JSON 并计算 \(R_{ls},R_{pm}\)）。
  \item \textbf{根型分裂 LaTeX 生成物：}\texttt{\detokenize{sections/generated/tab_sp6_isotropy_root_split_ratio.tex}}。
  \item \textbf{根型分裂审计 JSON：}\texttt{\detokenize{artifacts/export/sp6_isotropy_root_split_ratio.json}}。
\end{itemize}

\paragraph{接口延拓：将外部时间缩放、折叠缺陷与耦合漂移归约到同一主标量变量}
本章在注 \ref{rem:lapse-modulated-tau-mix} 中已将外部时间重标定写成一个最小的外部时间缩放代理
\(
N:=\kappa_0/\kappa
\)
（其中 \(\kappa\) 是背景依赖的路由/逆搜索开销，单位为“每内部步的外部代价”）。
另一方面，折叠侧已经定义了纤维缺陷标量
\begin{equation}\label{eq:gut_kappa_m_defect}
\kappa_m(\pi):=\mathbb{E}_{\pi}[\log d_m(X)]
\qquad(\text{定义 \ref{def:pom-kappa}；其中 }d_m(x)=|\Fold_m^{-1}(x)|).
\end{equation}
因此可以引入一个统一的主标量变量（纯定义，不引入额外指称）：
\[
u:=\log\frac{\kappa}{\kappa_0}=-\log N.
\]
\emph{接口级}假设是：在带空间索引的读出（例如 Hilbert 地址化窗口）下，局部开销场 \(u(x)\) 与某个局部折叠统计的缺陷量同构，
例如以局部分布 \(\pi_{m,x}\in\Delta(X_m)\) 读出
\(
\kappa_m(\pi_{m,x})
\)
并设
\[
u(x)\ \approx\ c\,\kappa_m(\pi_{m,x})
\qquad(c>0\ \text{为标定常数}).
\]
一旦将任意有效耦合或有效常数的漂移写成主标量的函数（线性化形式）
\[
g_i^{-2}(x)=a_i+b_i\,u(x)+\cdots,
\]
就得到一个\textbf{可证伪的相关性}：
\[
\Delta g_i^{-2}\ \approx\ b_i\,\Delta u\ =\ -b_i\,\Delta\log N.
\]
进一步地，对任意两指标 \(i,j\)（且 \(b_j\neq 0\)），由同一线性化立即得到漂移比值不变量
\[
\frac{\Delta g_i^{-2}}{\Delta g_j^{-2}}\ \approx\ \frac{b_i}{b_j},
\]
从而“多耦合漂移”在该口径下被压缩为由单一潜变量 \(u\) 参数化的一条仿射共线族：所有 \((\Delta g_i^{-2})_i\) 在一阶近似上落在同一秩 \(1\) 方向上。
该相关性与比值不变量本身不依赖连续极限；它们只要求在同一数据/同一协议族上同时输出 \(N\) 与 \(\{g_i\}\) 的代理量并做联合检验。
\par
更完整的 overhead\(\to\)外部时间缩放\(\to\)线性化闭合，以及从地址化窗口与折叠退化统计重建空间代价标量变量的可执行协议，属于接口层的进一步构造；
本稿在此仅记录其在本章变量上的最小同构关系与可证伪相关性形式。

\begin{corollary}[纤维缺陷标量作为单潜变量：漂移的仿射共线性与比值不变量]\label{cor:gut-fiber-defect-single-latent-affine-collinearity}
沿用上段的主标量 \(u=-\log N\) 与一阶线性化 \(g_i^{-2}(x)=a_i+b_i\,u(x)+\cdots\)。
则对任意两观测点（或两组协议设定）\(x_1,x_2\) 的差分 \(\Delta(\cdot):=(\cdot)(x_2)-(\cdot)(x_1)\)，有
\[
\Delta g_i^{-2}\ \approx\ b_i\,\Delta u,
\qquad
\frac{\Delta g_i^{-2}}{\Delta g_j^{-2}}\ \approx\ \frac{b_i}{b_j}\ (b_j\neq 0),
\]
从而多耦合漂移在一阶近似上必落在同一秩 \(1\) 的共线方向上；该比值不变量可作为与连续极限无关的联合可证伪检验。
\leanverified{paper\_gut\_fiber\_defect\_affine\_collinearity\_seeds}
\end{corollary}
\begin{proof}
由线性化形式对 \(x\) 作差分立得第一式；对两指标取比即得第二式。证毕。
\end{proof}

\paragraph{有限检验模板：Hecke 刚性与 Sturm 界的有限素数骨架判定}
\begin{theorem}[审计化 Langlands 型识别的有限可判定性：Sturm 有界性 \(\times\) finite-layer closure \(\times\) counterterm strictification]\label{thm:audit-finite-decidability-sturm-finite-layer}
设某一候选“统一对象”的读出被组织为一组可比的系数/局部因子数据（例如模形式的 \(q\)-展开系数、Hecke 本征值 \(a_p\)，或等价的素数局部因子/特征多项式），并且其实现链允许同时输出 Dirichlet 扭曲通道与有限部漂移预算。
则在本文的可审计口径下，Langlands 型“识别/匹配”（是否属于给定的 Hecke 本征包/等价类）可归约为\textbf{有限}核验问题，其有限性来源可明确分解为三条：
\begin{enumerate}
  \item \textbf{（模形式端：Sturm 有界性）}在给定权 \(k\) 与水平 \(N\) 后，Sturm 界 \(B_{k,N}\) 给出有限截断判定：只需核验至 \(B_{k,N}\) 的有限个 Fourier/Hecke 数据即可判定等同性（从而把“无限条件”替换为有限素数骨架的检查）。
  \item \textbf{（Dirichlet 端：有限层闭包）}任何有限投影词/有限实现对 Dirichlet 轴束的依赖必在某个有限层 \(m_\ast\) 上闭合；因而“全角色扫描/最坏通道”可归约为对有限集合 \(\widehat{G_{m_\ast}}\) 的穷举核验（命题 \ref{prop:pom-finite-modulus-stability}，亦见推论 \ref{cor:pom-local-finite-dirichlet-action}）。
  \item \textbf{（接口严格化端：counterterm 消去有限部漂移）}中心扩张口径下的交换失败完全落在有限部/常数层，并可由纯常数 counterterm 门严格消去，从而把残差交换提升为严格相容性条件（命题 \ref{prop:pom-counterterm-anom-cancel}，亦见定理 \ref{thm:pom-strictification-splitting}）。
\end{enumerate}
因此，在既定的误差预算结构（有限样本/截断/strictification）下，上述识别问题不再是开放式的无穷验证，而是一个由 \((B_{k,N},m_\ast)\) 与有限部 counterterm 预算共同控制规模的有限判定问题。
\leanverified{paper\_audit\_finite\_decidability\_sturm\_finite\_layer}
\end{theorem}

\begin{remark}[总结性结论：lapse--Dirichlet--Adams--counterterm--Sturm 的闭合主线]\label{rem:audit-summary-lapse-dirichlet-adams-sturm}
本稿把“可审计计算”系统性地压缩为有限证书链：lapse（测度变换/外部时间缩放）可内化为加权周期读出并给出显式有限样本上界（定理 \ref{thm:op-algebra-lapse-weighted-wish-koksma}）；Dirichlet 扭曲通道在有限层闭包下局部有限并可穷举核验（命题 \ref{prop:pom-finite-modulus-stability} 与推论 \ref{cor:pom-local-finite-dirichlet-action}）；Adams--Dirichlet 的算术伸缩门把 \((\chi,u)\) 的幂半群作用统一为重整化单子（定义 \ref{def:witt-adams-dirichlet-renorm-gate} 与命题 \ref{prop:witt-renorm-monoid}）；异常通过中心扩张分裂与 counterterm strictification 被严格化为常数层可控漂移（定理 \ref{thm:pom-strictification-splitting}）。
在此基础上，Sturm 有界性把 Hecke 刚性落实为有限素数骨架判定（定理 \ref{thm:audit-finite-decidability-sturm-finite-layer}），从而形成一条无需手工编号、可逐条引用的闭合结构主线。
现有脚本 \texttt{\detokenize{scripts/exp_hecke_e4_prime_audit.py}} 即提供该模板的最小可复核实例：在零参数基准情形下，以有限素数点的 Hecke 系数表/JSON 输出作为快速否定检验。
\end{remark}

\begin{corollary}[Sturm--Hecke 有界性下的有限素数骨架否定判据：三素数即可排除标量 \texorpdfstring{$X(1)$}{X(1)} 权 \texorpdfstring{$4$}{4} 情形]\label{cor:hecke-e4-three-prime-refutation}
\leanverified{paper\_hecke\_e4\_three\_prime\_refutation}
若顶层统一对象被声明为 \emph{level-1（\(X(1)\)）标量权 \(4\) Hecke-闭合}的模形式，
则该空间一维且由 Eisenstein 级数 \(E_4\) 张成，因此其素数 Hecke 系数必须满足
\[
\boxed{\ a_p=240(1+p^3)\qquad(p\ \text{素})\ }.
\]
从而在最小否定口径下，仅需在 \(p\in\{2,3,5\}\) 三个素数点比对系数即可形成快速否定检验：
任一显著偏离即推出“顶层统一不可能仍停留在 \(X(1)\) 标量一维 + Hecke 刚性”层级，必须进入向量值/同余层提升分支。
\end{corollary}

\begin{proposition}[统一反例的类型图核归约]\label{prop:distill-group-unification-counterexample-kernel}
设 \(\Gamma\) 为有限类型邻接图，\(H\) 为所有保持给定审计读出数据的有限重标号群，并设存在同态
\[
\rho:H\longrightarrow \mathrm{Aut}(\Gamma).
\]
若 \(\mathrm{Aut}(\Gamma)=\{1\}\)，则任何可由类型重标号实现的统一反例均落在 \(\ker\rho=H\) 的类型图平凡作用部分；换言之，它不能改变任何类型邻接关系。

进一步，令 \(G\) 为紧连通半单 Lie 群，\(Z(G)\) 为其有限中心。
若连续包络的可观测作用由连续同态
\[
\kappa:G\longrightarrow H
\]
给出，则 \(\kappa\) 必为平凡同态。
若此外存在有限维实表示 \(V\)、向量 \(v_0\in V\) 与有限多项式审计族 \(\mathcal I\subset\mathbb R[V]\)，使映射
\[
G/Z(G)\longrightarrow \mathbb R^{\mathcal I},\qquad
 gZ(G)\longmapsto (I(gv_0))_{I\in\mathcal I}
\]
为单射，则任何同时保持类型图与全部审计值 \((I(v_0))_{I\in\mathcal I}\) 的连续包络元素必属于 \(Z(G)\)。
因此，在这些假设下，不存在正维的最坏统一反例；终端未消除自由度至多为有限中心商，即低复杂度 slab 型残余。
\end{proposition}
\begin{proof}
第一句由定义立即得到：若 \(\mathrm{Aut}(\Gamma)=\{1\}\)，则 \(\rho(H)\) 只能是平凡群，故任何审计保持重标号都不能在类型邻接图上产生非平凡作用。

因为 \(G\) 连通而 \(H\) 有限且离散，连续像 \(\kappa(G)\) 既连通又有限，故只能为单点；于是 \(\kappa\) 平凡。
最后，若 \(g\in G\) 保持全部审计值，则
\[
(I(gv_0))_{I\in\mathcal I}=(I(v_0))_{I\in\mathcal I}.
\]
由所给单射性，\(gZ(G)=Z(G)\)，即 \(g\in Z(G)\)。
半单紧连通群的中心有限，故剩余歧义不是正维连续自由度，而只能是有限中心商。
这正给出统一反例的类型图核归约与终端 slab 分类。证毕。
\end{proof}

\par
我们将最小否定判据素数点 \(p\in\{2,3,5\}\) 的数值表与 JSON 封装为可复现输出（作为“零拟合必要条件”的快速否定检验）：
\begin{itemize}
  \item \textbf{生成脚本：}\texttt{\detokenize{scripts/exp_hecke_e4_prime_audit.py}}。
  \item \textbf{LaTeX 生成物：}\texttt{\detokenize{sections/generated/tab_hecke_e4_prime_audit_p235.tex}}。
  \item \textbf{审计 JSON：}\texttt{\detokenize{artifacts/export/hecke_e4_prime_audit_p235.json}}。
\end{itemize}
